\documentclass[../../main.tex]{subfiles}

\begin{document}

\section{Ascoli 定理}\label{section:ch07-ascoli}

我们现在来证明 Ascoli 定理的一般形式, 它刻画了关于紧致收敛拓扑 $\mathcal{C}(X,Y)$ 的紧致子空间. 然而, 这个证明涉及函数空间的三种标准拓扑, 即点态收敛拓扑、紧致收敛拓扑和一致拓扑.

\begin{theorem}[Ascoli 定理]\label{theorem:ch07-ascoli-general}
设 $X$ 是一个空间, $(Y,d)$ 是一个度量空间. 赋予 $\mathcal{C}(X,Y)$ 紧致收敛拓扑. 设 $\mathcal{F}$ 是 $\mathcal{C}(X,Y)$ 的一个子集.
\begin{enumerate}[(a)]
\item 若 $\mathcal{F}$ 关于 $d$ 是等度连续的, 并且集合
\begin{align*}
\mathcal{F}_a=\{f(a)\mid f\in\mathcal{F}\}
\end{align*}
对于每一个 $a\in X$ 有紧致闭包, 则 $\mathcal{F}$ 包含于 $\mathcal{C}(X,Y)$ 的一个紧致子空间中.
\item 若 $X$ 是一个局部紧致的 Hausdorff 空间, 则逆命题成立.
\end{enumerate}
\end{theorem}

\begin{proof}
(a) 的证明. 赋予 $Y^X$ 积拓扑, 它与点态收敛拓扑是一致的. 从而 $Y^X$ 是一个 Hausdorff 空间. 空间 $\mathcal{C}(X,Y)$ 有紧致收敛拓扑, 它不是 $Y^X$ 的子空间. 设 $\mathcal{G}$ 是 $\mathcal{F}$ 在 $Y^X$ 中的闭包.

第一步. 证明 $\mathcal{G}$ 是 $Y^X$ 的一个紧致子空间. 给定 $a\in X$, 令 $C_a$ 表示 $\mathcal{F}_a$ 在 $Y$ 中的闭包. 根据假设条件可见, $C_a$ 是 $Y$ 的一个紧致子空间. 集合 $\mathcal{F}$ 包含于积空间
\begin{align*}
\prod_{a\in X} C_a,
\end{align*}
因为根据定义, 这个积空间由满足以下条件的所有函数 $f\colon X\to Y$ 构成: 对于所有 $a$, 有 $f(a)\in C_a$. 根据 Tychonoff 定理, 这个积空间是紧致的, 它是 $Y^X$ 的一个闭子空间. 因为 $\mathcal{G}$ 是 $\mathcal{F}$ 在 $Y^X$ 中的闭包, 所以 $\mathcal{G}$ 包含于 $\prod C_a$ 中, 并且是闭的. 从而 $\mathcal{G}$ 是紧致的.

第二步. 证明 $\mathcal{G}$ 中的每一个函数都是连续的, 并且 $\mathcal{G}$ 关于 $d$ 是等度连续的.

给定 $x_0\in X$ 和 $\varepsilon>0$, 选取 $x_0$ 的一个邻域 $U$, 使得
\begin{align*}
d(f(x),f(x_0))<\varepsilon/3\quad\text{对于所有 }f\in\mathcal{F}\text{ 和所有 }x\in U.
\end{align*}
下面将证明对于 $g\in\mathcal{G}$ 和 $x\in U$, 有 $d(g(x),g(x_0))<\varepsilon$. 从而 $\mathcal{G}$ 是等度连续的.

设 $g\in\mathcal{G}$, $x$ 是 $U$ 的一个点. 定义 $V_x$ 是 $Y^X$ 中所有元素 $h$ 构成的一个开集, 其中 $h$ 满足条件
\begin{align*}
d(h(x),g(x))<\varepsilon/3\text{ 和 }d(h(x_0),g(x_0))<\varepsilon/3.
\end{align*}
因为 $g$ 属于 $\mathcal{F}$ 的闭包, 所以 $g$ 的邻域 $V_x$ 必然包含 $\mathcal{F}$ 的一个元素 $f$, 对于上述两个不等式使用三角不等式, 可见 $d(g(x),g(x_0))<\varepsilon$.

第三步. 证明 $Y^X$ 上的积拓扑和 $\mathcal{C}(X,Y)$ 上的紧致收敛拓扑在子集 $\mathcal{G}$ 上是一致的.

一般地, 紧致收敛拓扑要比积拓扑细. 我们证明对于子集 $\mathcal{G}$, 相反的蕴涵关系也成立. 令 $g$ 为 $\mathcal{G}$ 的一个元素, 令 $B_C(g,\varepsilon)$ 是 $Y^X$ 上紧致收敛拓扑的包含 $g$ 的一个基元素, 要找到 $Y^X$ 上点态收敛拓扑的一个包含 $g$ 的基元素 $B$, 使得
\begin{align*}
[B\cap\mathcal{G}]\subset[B_C(g,\varepsilon)\cap\mathcal{G}].
\end{align*}
应用 $\mathcal{G}$ 的等度连续性和 $C$ 的紧致性, 可以选出 $X$ 的有限多个开集 $U_1,\cdots,U_n$ 覆盖 $C$, 要求 $U_1,\cdots,U_n$ 分别包含点 $x_1,\cdots,x_n$, 并且对于每一个 $i$, 有
\begin{align*}
d(g(x),g(x_i))<\varepsilon/3
\end{align*}
其中 $x\in U_i$, $g\in\mathcal{G}$. 定义 $Y^X$ 的基元素 $B$ 为
\begin{align*}
B=\{h\mid h\in Y^X\text{ 和 }d(h(x_i),g(x_i))<\varepsilon/3,i=1,\cdots,n\}.
\end{align*}
下面证明若 $h$ 是 $B\cap\mathcal{G}$ 中的一个元素, 则 $h$ 便属于 $B_C(g,\varepsilon)$. 也就是证明, 对于 $x\in C$ 有 $d(h(x),g(x))<\varepsilon$. 给定 $x\in C$, 选取 $i$ 使得 $x\in U_i$. 由于 $x\in U_i$ 和 $g,h\in\mathcal{G}$, 所以有
\begin{align*}
d(h(x),h(x_i))<\varepsilon/3\text{ 和 }d(g(x),g(x_i))<\varepsilon/3.
\end{align*}
由于 $h\in B$, 所以
\begin{align*}
d(h(x_i),g(x_i))<\varepsilon/3.
\end{align*}
从而根据三角不等式可见 $d(h(x),g(x))<\varepsilon$.

第四步. 完成证明. 集合 $\mathcal{G}$ 包含 $\mathcal{F}$ 并且包含于 $\mathcal{C}(X,Y)$, 并且它是 $Y^X$ 关于积拓扑的一个紧致子空间. 根据刚刚证明的结论, 它也是 $\mathcal{C}(X,Y)$ 关于紧致收敛拓扑的一个紧致子空间.

(b) 的证明. 设 $\mathcal{H}$ 是 $\mathcal{C}(X,Y)$ 中包含 $\mathcal{F}$ 的一个紧致子空间. 下面证明 $\mathcal{H}$ 是等度连续的, 并且对于每一个 $a\in X$, $\mathcal{H}_a$ 是紧致的. 由此可见 $\mathcal{F}$ 是等度连续的 (由于 $\mathcal{F}\subset\mathcal{H}$), 以及 $\mathcal{F}_a$ 包含于 $Y$ 的紧致子空间 $\mathcal{H}_a$, 从而 $\overline{\mathcal{F}_a}$ 是紧致的.

为了证明 $\mathcal{H}_a$ 是紧致的, 考虑以下映射的复合:
\begin{align*}
j\colon\mathcal{C}(X,Y)\longrightarrow X\times\mathcal{C}(X,Y),
\end{align*}
其中 $j(f)=a\times f$, 以及赋值映射
\begin{align*}
e\colon X\times\mathcal{C}(X,Y)\longrightarrow Y,
\end{align*}
其中 $e(x\times f)=f(x)$. 显然映射 $j$ 是连续的, 根据\cref{theorem:ch07-compact-open-equals-compact-convergence}和\cref{theorem:ch07-evaluation-continuous}可见, 映射 $e$ 也是连续的. 复合映射 $e\circ j$ 将 $\mathcal{H}$ 映到 $\mathcal{H}_a$ 中, 因为 $\mathcal{H}$ 是紧致的, 所以 $\mathcal{H}_a$ 也是紧致的.

现在我们来证明 $\mathcal{H}$ 在 $a$ 点关于度量 $d$ 是等度连续的. 设 $A$ 是 $X$ 的一个紧致子空间, 包含 $a$ 的一个邻域. 下面只需证明 $\mathcal{C}(A,Y)$ 的子集
\begin{align*}
\mathcal{R}=\{f\mid A;f\in\mathcal{H}\}
\end{align*}
在 $a$ 点处是等度连续的.

赋予 $\mathcal{C}(A,Y)$ 紧致收敛拓扑, 下面证明限制映射
\begin{align*}
r\colon\mathcal{C}(X,Y)\longrightarrow\mathcal{C}(A,Y)
\end{align*}
是连续的. 设 $f$ 是 $\mathcal{C}(X,Y)$ 中的一个元素, $B=B_C(f\mid A,\varepsilon)$ 是 $\mathcal{C}(X,Y)$ 的包含 $f\mid A$ 的一个基元素, 其中 $C$ 是 $A$ 的一个紧致子空间. 那么 $C$ 是 $X$ 的一个紧致子空间, 并且 $r$ 将 $\mathcal{C}(X,Y)$ 中 $f$ 的邻域 $B_C(f,\varepsilon)$ 映到 $B$ 中.

映射 $r$ 将 $\mathcal{H}$ 映到 $\mathcal{R}$ 上. 因为 $\mathcal{H}$ 是紧致的, 所以 $\mathcal{R}$ 也是紧致的. $\mathcal{R}$ 是 $\mathcal{C}(A,Y)$ 的一个子空间, 因为 $A$ 是紧致的, 紧致收敛拓扑和一致收敛拓扑在 $\mathcal{C}(A,Y)$ 上是一致的. 根据\cref{theorem:ch07-compact-complete-totally}可见, $\mathcal{R}$ 关于 $\mathcal{C}(A,Y)$ 的一致度量是完全有界的. 从而根据\cref{lemma:ch07-totally-bounded-equicontinuous}可见, $\mathcal{R}$ 相对于度量 $d$ 是等度连续的.
\end{proof}

至于 Ascoli 定理的更为一般的形式可以在 [K] 或 [Wd] 中找到. 在那里并没有要求 $Y$ 是一个度量空间, 仅仅假定它有所谓的一致结构. 这种一致结构是度量概念的一种推广.

Ascoli 定理在分析中有很多应用, 这已超出本书的范围. 可以在 [K-F] 中找到一些这样的应用.

\subsection*{习题}

\begin{enumerate}
\item 以下所列举的 $\mathcal{C}(\mathbb{R},\mathbb{R})$ 的一些子集中哪些是点态有界的? 哪些是等度连续的?
\begin{enumerate}[(a)]
\item 族 $\{f_n\}$, 其中 $f_n(x)=x+\sin nx$.
\item 族 $\{g_n\}$, 其中 $g_n(x)=n+\sin x$.
\item 族 $\{h_n\}$, 其中 $h_n(x)=|x|^{1/n}$.
\item 族 $\{k_n\}$, 其中 $k_n(x)=n\sin(x/n)$.
\end{enumerate}
\item 证明:
\begin{quote}
\textbf{定理.} 若 $X$ 是一个局部紧致的 Hausdorff 空间, 则 $\mathcal{C}(X,\mathbb{R}^n)$ 关于紧致收敛拓扑的子空间 $\mathcal{F}$ 有紧致闭包, 当且仅当在 $\mathbb{R}^n$ 的任意一个标准度量下, $\mathcal{F}$ 是点态收敛的并且等度连续的.
\end{quote}
\item 证明: 当 $X$ 是一个 Hausdorff 空间时, Ascoli 定理的一般形式蕴涵经典形式 (\cref{theorem:ch07-ascoli-classical}).
\item 证明:
\begin{quote}
\textbf{定理(Arzelà 定理, 一般形式).} 设 $X$ 是一个 $\sigma$-紧致的 Hausdorff 空间, $f_n\colon X\to\mathbb{R}^k$ 是一个函数序列. 若 $\{f_n\}$ 是点态有界的和等度连续的, 则序列 $f_n$ 存在一个子序列, 它关于紧致收敛拓扑收敛到一个连续函数.
\end{quote}
[提示: 证明 $\mathcal{C}(X,\mathbb{R}^k)$ 是第一可数的.]
\item 设 $(Y,d)$ 是一个度量空间, $f_n\colon X\to Y$ 是一个连续函数序列, $f\colon X\to Y$ 是一个函数 (不必是连续的). 设关于点态收敛拓扑 $f_n$ 收敛到 $f$. 证明, 若 $\{f_n\}$ 是等度连续的, 则 $f$ 是连续的, 并且 $f_n$ 关于紧致收敛拓扑收敛到 $f$.
\end{enumerate}

\end{document}